\subsection*{§. 192.\\ Cayley 度量。}

刚得到的结果引我们讨论\emph{Cayley 度量}。这种度量是 Cayley 在 Beltrami 对该问题的研究以前建立的；下面将看到，它在实质上等价于常曲率空间的度量，并对这些空间的几何学给出一个十分清楚而简单的图像。

以任意二次曲面代替 Beltrami 超球面
\[
 x_1^2+x_2^2+\cdots+x_n^2=1
\]
并设其方程
\[
 f(x_1,x_2,\ldots,x_n)=0
\]
具有实系数；此外再附加一个条件（其意义马上就会看到）：\emph{该二次曲面上不存在实直线。}任取空间中的两点 $P,P'$，把它们的距离定义为交比
\[
 (PP'AB),
\]
的对数再乘以一个常数，其中 $A,B$ 是直线 $PP'$ 与二次曲面的两个交点。

若限于这样的空间区域：从其中各点不能向二次曲面作实切线，那么显然，如此定义的距离仅在两点 $P,P'$ 重合时为零；于是按照基本约定，$ds^2$ 将由一个\emph{正定}二次微分形式表示。正因如此，必须假定二次曲面不含实直线；否则，从空间的每一点都能向二次曲面作实切线。

再注意，基本二次曲面上的点在 Cayley 度量中表示距离为无穷大的点；因此，这个二次曲面也称为\emph{绝对形}。

容易\emph{先验地}看出，Cayley 度量就是常曲率空间的度量。事实上，众所周知，$S_n$ 中的任意二次曲面容许一个含 $\frac{n(n+1)}2$ 个参数、把二次曲面变为自身的连续射影变换群；而射影变换保持交比，也就保持 Cayley 度量的距离，所以对于具有相应 $ds^2$ 的空间，这是一个连续运动群，因而空间曲率为常数（§. 160）。

还要指出，在 Cayley 度量中，测地线由直线表示。事实上，任取一条不与绝对形相切的直线（一个 $S_1$）及其极线性空间 $S_{n-2}$；后者与 $S_1$ 没有公共点。以 $S_1,S_{n-2}$ 为不动点空间的调和轴同调，是二次曲面到自身的一个射影变换；在这个变换中，$S_1$ 的所有点都保持不动，而 $S_1$ 邻域内没有别的点保持不动\srcfn{1}{除 $S_1$ 的点外，空间中保持不动的只有其极空间 $S_{n-2}$ 的点。}。

因此，若 $A,B$ 是 $S_1$ 的任意两点，则连接它们的测地线也必须保持不动，故它与 $S_1$ 重合。

现在引入齐次坐标
\[
 x_0,x_1,x_2,\ldots,x_n,
\]
把二次曲面的方程写成
\[
 \mathop{\mathrm{S}}_{r,s} a_{rs}x_r x_s=0.
\]
由于它不应含有直线，按已知方法用实代换把左端的形式 $\Omega$ 化为平方和时，各平方项的系数必须全同号，或者除一个以外全同号。因此，不失一般性，经过射影变换后可以令
\[
 \Omega_{xx}=x_0^2+x_1^2+x_2^2+\cdots+x_n^2,
\]
或者
\[
 \Omega_{xx}=-x_0^2+x_1^2+x_2^2+\cdots+x_n^2.
\]
第二种情形显然使我们回到伪球空间的 Beltrami 测地表示（§. 191），即双曲度量；第一种情形则给出现在要讨论的椭圆度量（具有虚绝对形）。

